\documentclass[11pt, a4paper]{article}

% Préambule : chargement des packages et configuration
\usepackage[utf8]{inputenc}   % Encodage des caractères (gère les accents)
\usepackage[T1]{fontenc}      % Encodage des polices
\usepackage[french]{babel}    % Langue du document (ici en français)
\usepackage{amsmath, amssymb} % Pour les mathématiques avancées

\title{Mon Premier Document en LaTeX}
\author{Votre Nom}
\date{\today}

% Corps du document
\begin{document}

\maketitle

\section{Introduction}
Ceci est un paragraphe rédigé en LaTeX. La mise en page et la gestion des espaces sont automatiques.

\section{Mathématiques}
LaTeX excelle dans l'écriture de formules mathématiques. Par exemple, le théorème de Pythagore s'écrit :
\begin{equation}
    a^2 + b^2 = c^2
\end{equation}

\end{document}